\clearpage
\SourceLeaf{402}\pdfbookmark[1]{II. Tafeln für die Störungen der Ceres}{ceres-tables-letters}
\begin{center}
{\large BRIEFWECHSEL.\par}
\bigskip
[II. TAFELN FÜR DIE STÖRUNGEN DER CERES.]\par
\bigskip
\textsc{Gauss} an \textsc{Olbers}, Braunschweig, 3. December 1802.
\end{center}
\noindent\hspace*{\parindent}.\quad.\quad.\quad.\quad.\quad Ich berechne mir jetzt eine Tafel für die im \rlap{\kern.02em\raisebox{1.02em}{\rule{.2em}{.25pt}}}Nov. der M. C. a\rlap{\kern.35em\raisebox{1.02em}{\rule{.2em}{.25pt}}}ngezeigten Störungen der {\AstroSymbols ⚳}, wobei ich zum Theil eine Einrichtung mache, die von der sonst üblichen verschieden ist. Die Tafeln für die bloss von \(\text{\AstroSymbols ⚳}-\text{\AstroSymbols ♃}\) abhängenden Störungen der Länge und des Radius vector und für die Breite sind schon fertig. Ich schreibe Ihnen darüber nächstens ausführlicher.\quad.\quad.\quad.\quad.\quad.

\begin{center}
\textsc{Gauss} an \textsc{Olbers}, Braunschweig, 21. December 1802.
\end{center}
\noindent\hspace*{\parindent}.\quad.\quad.\quad.\quad.\quad Meine, jetzt vollendeten, Tafeln\refstepcounter{footnote}\hyperlink{ceresII-note}{[*)]}\begingroup
\makeatletter\long\def\@makefntext#1{\noindent[*) #1}\makeatother
\footnotetext{\hypertarget{ceresII-note}{}Band VI, Seite 235—243.]}
\endgroup{} der Störungen der {\AstroSymbols ⚳} vom {\AstroSymbols ♃} werde ich noch in dieser Woche an \textsc{Zach} schicken; sind sie für die M. C. noch zu weitläufig, so werde ich Ihnen eine Abschrift davon besorgen. Ihnen kann ich die Einrichtung mit ein paar Worten beschreiben. Bei der Länge machen erstlich alle Gleichungen, deren Argument ein Vielfaches von \(\text{\AstroSymbols ⚳}-\text{\AstroSymbols ♃}\) ist, Eine Tafel; die Summe aller übrigen, von der einfachen Excentricität abhängigen, lässt sich durch \(A\sin(B-\text{\AstroSymbols ⚳})\) ausdrücken, so dass \(A\) und \(B\) bloss von \(\text{\AstroSymbols ⚳}-\text{\AstroSymbols ♃}\) Functionen sind; diese Grössen \(A\) und \(B\) gibt, für alle Werthe von \(\text{\AstroSymbols ⚳}-\text{\AstroSymbols ♃}\) von Grad zu Grad, eine zweite Tafel. Eben so kommen 2 Tafeln für den Radius vector und Eine für die Breite, zusammen 5 Tafeln, da man sonst für alle die von mir gebrauchten Gleichungen 40 Tafeln oder doch an 30 brauchen \SourceLeaf{403}würde, wenn man die wegliesse, die nicht auf \(2''\) gehen. Mir wenigstens ist es \textls[180]{weit} leichter, diese leichten trigonometrischen Rechnungen zu machen, als so viele Argumente zu bilden, in so viele Tafeln einzugehen, so viele Additionen zu machen. Man könnte auch leicht den Gebrauch meiner Tafel noch mehr erleichtern, wenn man z. B. statt \(A\) gleich \(\log A\) aufnähme; oder auch, wenn man nach Gefallen eine beständige Grösse \(=2M\) annähme, die nicht kleiner, als irgend ein Werth von \(A\) wäre, \(\frac{A}{M}=\cos C\) machte: hiedurch würde die Summe der von der Excentricität abhängigen Gleichungen
\[
=M\sin(B-C-\text{\AstroSymbols ⚳})+M\sin(B+C-\text{\AstroSymbols ⚳}).
\]
Man nähme dann in die Tafel, für jeden Werth von \(\text{\AstroSymbols ⚳}-\text{\AstroSymbols ♃}\) die Werthe von \(B-C\) und von \(B+C\) auf und fügte noch eine Tafel hinzu, die \(M\sin\varphi\) für alle Werthe von \(\varphi\) enthielte. — In meiner Tafel habe ich übrigens nicht die Winkel \(B\), sondern für:
\[
\begin{array}{@{}l@{\quad}l@{}}
\text{die Länge}\quad\text{.\ .\ .\ .}&B-2(\text{\AstroSymbols ⚳}-\text{\AstroSymbols ♃})\\[3pt]
\left.\begin{array}{@{}l@{}}\text{den R. V.}\\\text{Breite}\end{array}\right\}\ \text{.\ .\ .}&B-3(\text{\AstroSymbols ⚳}-\text{\AstroSymbols ♃})
\end{array}
\]
angesetzt oder der Summe der Gleichungen die Form
\[
\begin{array}{@{}l@{\quad}l@{}}
\text{für die Länge}\ \text{.\ .\ .\ .\ .\ .\ .\ .\ .\ .\ .\ .\ .\ .\ .}&A\sin(B'+\text{\AstroSymbols ⚳}-2\text{\AstroSymbols ♃})\\[4pt]
\text{für den R. V. und die Br.}\ \text{.\ .\ .\ .}&A\sin(B'+2\text{\AstroSymbols ⚳}-3\text{\AstroSymbols ♃})
\end{array}
\]
gegeben.

Dieses \(B'\) bleibt so immer zwischen bestimmten Grenzen, wenn \(\text{\AstroSymbols ⚳}-\text{\AstroSymbols ♃}\) alle Werthe durchläuft, da hingegen \(B\) einen oder zwei Umläufe machen würde während einer Periode von \(\text{\AstroSymbols ⚳}-\text{\AstroSymbols ♃}\).\quad.\quad.\quad.\quad.\quad.

\begin{center}
\textsc{Gauss} an \textsc{Olbers}, Braunschweig, 4. Januar 1803.
\end{center}
\noindent\hspace*{\parindent}.\quad.\quad.\quad.\quad.\quad Nach meinen vorläufigen Tafeln der {\AstroSymbols ⚳}-Störungen .\quad.\quad.\quad.\quad habe ich schon häufige Rechnungen gemacht, und ihre vorzügliche \textls[180]{Bequemlichkeit} bewährt gefunden. Aber mit den VIII. im Nov. der M. C. abgedruckten Elementen bin ich noch nicht recht zufrieden. Ich hatte dazu lediglich Meridianbeobachtungen gebraucht, die dadurch sehr gut, und viel besser als durch die\,\raisebox{.65ex}{.} VII. dargestellt wurden. Aber jetzt, nachdem die ohne Tafeln so \SourceLeaf{404}sehr beschwerlichen Berechnungen der Störungen kein Hinderniss mehr in den Weg legten, habe ich auch die letzten \textsc{Oriani}schen Beobachtungen damit verglichen, und zu meiner Verwunderung den Fehler merklich grösser gefunden, als nach den VII. Elementen. Die Ascensionen stimmen alle ziemlich gut, aber die Declinationen sind mit vieler Harmonie etwa \(40''\) zu klein, so dass ich den Längenfehler für den 5. Aug. \(+18''\), den Breitenfehler \(-37''\) im Mittel ansetze; der Fehler der VII. Elemente ist nur etwa halb so gross. Mich dünkt, dies bestätigt eine schon ehemals von mir geäusserte Meinung, dass die Unsicherheit, die eine Folge der noch zu kurzen Dauer der Beobachtungen ist, bis jetzt auf die Voraussagung des künftigen Laufs nachtheiliger wirken kann, als die Vernachlässigung der Störungen. Ich will deswegen erst noch eine Politur vornehmen und die Elemente allen vorhandenen Beobachtungen möglichst genau anpassen. Fast alle Rechnungen, die ich dazu machen muss, müsste ich ohnehin in diesem Jahre 1803 zur Benutzung der künftigen Beobachtungen doch auch machen; die Arbeit ist also nicht verloren, ob sie gleich zur Wiederauffindung am östlichen Himmel eben nicht nöthig ist; dazu werden, denke ich, die von \textsc{Triesnecker} und \textsc{Bode} berechneten Ephemeriden hinreichend genau sein. — Was übrigens die Erreichung einer grossen Präcision in den Elementen noch hindert, ist hauptsächlich die kurze Dauer der \textsc{Piazzi}schen Beobachtungen von 1801. Zu einer genauen Bestimmung, wenn so wie hier die mittlere Bewegung nicht aus sehr entfernten Beobachtungen abgeleitet werden kann, ist es \textls[180]{wesentlich}, dass man 4 Beobachtungen \(A,B,C,D\) habe, die alle weit genug von einander liegen, auch wo möglich das Intervall \(B,C\) viel grösser sei, als \(A,B\) und\,\raisebox{.8ex}{.} \(C,D\). Dazu kann man nicht umhin, dass man für \(A\) und \(B\) die ersten und letzten \textsc{Piazzi}schen Beobachtungen nehme, die nur 41 Tage von einander abstehen. Man kann sich daher einen viel grössern Grad von Genauigkeit versprechen, wenn man von \textsc{Piazzi}s Beobachtungen nur noch Eine für \(A\) nöthig hat (die man natürlich durch Vergleichung mit den benachbarten genauer macht, als Eine einzelne ist), also 1803
\[\textstyle
\text{wo man }\left\{\begin{array}{c}A\\B\\C\\D\end{array}\right\}
\text{ aus }\left\{\begin{array}{ll}
1801&\\
1802&\text{(oder \textsc{Zach}s Beob. vom 7. Decbr. 1801)}\\
1802&\\
1803&
\end{array}\right.
\]
\SourceLeaf{405}nehmen kann. Aber eine recht grosse Genauigkeit darf man erst nach der {\AstroSymbols ☍} von 1804 hoffen, wenn man \(A,B,C,D\) aus so viel verschiedenen Jahren nehmen kann, und 1806, wenn man 4 Oppositionen benutzen kann, wird man hoffentlich die Bahn der {\AstroSymbols ⚳}, etwa die mittlere Bewegung ausgenommen, fast eben so genau kennen, als die der übrigen Planeten. Mit der {\AstroSymbols ⚴} wird alles Ein Jahr länger dauern. — Finden wir aber, woran wir noch nicht verzweifeln wollen, einen von beiden in einer ältern Beobachtung wieder, so wird freilich alles viel geschwinder gehen.\quad.\quad.\quad.\quad.\quad.

\begin{center}
\textsc{Gauss} an \textsc{Olbers}, Braunschweig, 1. März 1803.
\end{center}
\noindent\hspace*{\parindent}.\quad.\quad.\quad.\quad.\quad Meine Tafeln für die Störungen der {\AstroSymbols ⚳} hätte ich Ihnen längst geschickt, wenn nicht von \textsc{Zach} mir geschrieben hätte, dass alles im Februarhefte der M. C. erscheinen würde. Da in diesem der Platz nicht zugereicht haben wird, so erhalten Sie dieselben in meinem nächsten Briefe. Ich habe mir auch die übrigen für die Bewegung der {\AstroSymbols ⚳} nöthigen Tafeln leni calamo nur auf Minuten berechnet, eigentlich nur zu meinem Gebrauch, um die ganz neuerlich erhaltene Histoire Céleste durchzumustern; ich habe sie bereits durchblättert und einige flüchtige Versuche angestellt; ich verzweifle aber fast, die {\AstroSymbols ⚳} darin zu finden, ob ich gleich mehrere Tage angetroffen habe, wo der Unterschied der Declination nur einige Grade betrug. Auch diese Tafel schicke ich Ihnen nächstens. Beide Tafeln sind indess nur provisorisch; nach dem Schluss der diesjährigen Beobachtungen werden wir weit genauere Elemente haben, und dann wird es schon die Mühe belohnen, die Störungen genauer und vollständiger zu berechnen und aufzunehmen.\quad.\quad.\quad.\quad.\quad.

\begin{center}
\textsc{Gauss} an \textsc{Olbers}, Braunschweig, 8. April 1803.
\end{center}
\noindent\hspace*{\parindent}.\quad.\quad.\quad.\quad.\quad Hiebei schicke ich Ihnen die so lange versprochenen kleinen {\AstroSymbols ⚳}- und {\AstroSymbols ⚴}-Tafeln; die eine noch rückständige für die {\AstroSymbols ⚴} erhalten Sie nächstens nach.\quad.\quad.\quad.\quad.\quad.

Erst in gegenwärtigem Jahre werde ich zur fernern Verbesserung der {\AstroSymbols ⚳}-Elemente diejenige Methode anwenden können, die dann immerfort zur weitern Ausfeilung brauchbar sein wird; bisher war es dazu zu früh. Da diese \SourceLeaf{406}Methode meiner Meinung nach kein unwichtiger Theil des Ganzen ist, denke ich an die Bearbeitung des Ganzen dann zu gehen, wenn ich diesen Theil erst \textls[180]{praktisch} geprüft und, wie das zu gehen pflegt, ihn praktisch vervollkommnet habe.\quad.\quad.\quad.\quad.\quad.

\begin{center}
\textsc{Gauss} an \textsc{Olbers}, Braunschweig, 25. Januar 1805.
\end{center}
\noindent\hspace*{\parindent}.\quad.\quad.\quad.\quad.\quad Meine neuen Elemente der Ceres .\quad.\quad.\quad.\quad.\quad sind folgende:
\begin{center}
\begingroup\small\setlength{\tabcolsep}{5pt}\renewcommand{\arraystretch}{1.25}
\begin{tabular}{r r r r}
 & Epoche & Aphelium & Excentricität\\
1801 & \(77^0\ 17'\ \ 0''{,}1\) & \(326^0\ 19'\ 59''\) & \(0{,}078\ 4929\)\\
1802 & \(155\ \ 27\ \ 34{,}2\) & \(22\ \ \ 0\) & \(4871\)\\
1803 & \(233\ \ 38\ \ \ 8{,}3\) & \(24\ \ \ 2\) & \(4814\)\\
1804 & \(312\ \ \ 1\ \ 33{,}5\) & \(26\ \ \ 3\) & \(4757\)\\
1805 & \(30\ \ 12\ \ \ 7{,}7\) & \(28\ \ \ 4\) & \(4700\)
\end{tabular}\endgroup
\end{center}
\begin{center}
\begingroup\footnotesize\setlength{\tabcolsep}{2pt}\renewcommand{\arraystretch}{1.25}
\begin{tabular}{l@{\quad}r}
Logarithm der halben grossen Axe .\ .\ .\ .\ . & \(0{,}442\ 0004\)\\
Tägliche tropische Bewegung .\ .\ .\ .\ .\ .\ .\ .\ .\ . & \(771''{,}0524\)\\
\(\AscendingNode\) 1801 .\ .\ .\ .\ .\ .\ .\ .\ .\ .\ .\ .\ .\ .\ .\ .\ .\ .\ .\ .\ .\ .\ .\ .\ .\ .\ .\ . & \(80^0\ 54'\ 46''\)\\
Neigung 1801 .\ .\ .\ .\ .\ .\ .\ .\ .\ .\ .\ .\ .\ .\ .\ .\ .\ .\ .\ .\ .\ .\ . & \(10^0\ 38'\ 13''\).
\end{tabular}\endgroup
\end{center}

Diese Elemente, mit denen noch meine alten Perturbationstafeln verbunden werden müssen, sind auf die 3 {\AstroSymbols ☍} und \textsc{Piazzis} Beobachtungen von 1801 gegründet. Inzwischen lassen sich die \textls[180]{Breiten} in den 4 Jahren nicht mit der gefundenen Säcularänderung vom \(\AscendingNode\) und Inclination darstellen (das erstere bemerkte ich schon voriges Jahr); und ich muss, um die Beobachtungen mit meinen Störungstafeln zu vereinigen,
\begin{center}
\begin{tabular}{l l l}
dem Knoten & \(0''{,}243\) & tägl. Verrückung\\
der Neigung & \(0''{,}025\) & tägl. Abnahme
\end{tabular}
\end{center}
geben, beide viel grösser, als sie die Rechnung gibt. Ohne Zweifel liegt dies bloss an der Unvollständigkeit der periodischen Gleichungen für den Radius vector und die Breite, und wird sich künftig bei umständlicher Rechnung der Störungen schon ausweisen.
\par\medskip\begin{center}\rule{.22\linewidth}{.4pt}\end{center}
